\section{Polynomials}

\setcounter{subsection}{1}

\begin{exercise}[4a1]
    Suppose $w, z \in \cbb$.
	Verify the following equalities and inequalities.
    \begin{subexercises}
        \item $z + \bar{z} = 2 \Re z$
        \item $z - \bar{z} = 2(\Im z)i$
        \item $z\bar{z} = |z|^2$
        \item $\bar{w + z} = \bar{w} + \bar{z}$ and $\bar{wz} = \bar{w} \bar{z}$
        \item $\bar{\bar{z}} = z$
        \item $|\Re z| \leq |z|$ and $|\Im z| \leq |z|$
        \item $|\bar{z}| = |z|$
        \item $|wz| = |w| |z|$
    \end{subexercises}
\end{exercise}

\begin{sol}
    In this problem, let $z = a + bi$ and $w = c + di$ with $a, b, c, d \in \rbb$.
	Then
    \begin{subexercises}
        \item $z + \bar{z} = (a + bi) + (a - bi) = 2a = 2 \Re z$.
        \item $z - \bar{z} = (a + bi) - (a - bi) = 2bi = 2(\Im z)i$.
        \item $z\bar{z} = (a + bi)(a - bi) = a^2 + b^2 = |z|^2$.
        \item $\bar{w + z} = \bar{(c + di) + (a + bi)} = \bar{(a + c) + (b + d)i} = (a + c) - (b + d)i = (a - bi) + (c - di) = \bar{w} + \bar{z}$ and $\bar{wz} = \bar{(c + di)(a + bi)} = \bar{(ac - bd) + (ad + bc)i} = (ac - bd) - (ad + bc)i = (c - di)(a - bi) = \bar{w} \bar{z}$.
        \item $\bar{\bar{z}} = \bar{a - bi} = a + bi = z$.
        \item $|\Re z| = |a| \leq \sqrt{a^2 + b^2} = |z|$ and $|\Im z| = |b| \leq \sqrt{a^2 + b^2} = |z|$.
        \item $|\bar{z}| = |a - bi| = \sqrt{a^2 + b^2} = |z|$.
        \item $|wz| = |(c + di)(a + bi)| = |(ac - bd) + (ad + bc)i| = \sqrt{(ac - bd)^2 + (ad + bc)^2} = \sqrt{(a^2 + b^2)(c^2 + d^2)} = |a + bi| |c + di| = |z| |w|$. \qh
    \end{subexercises}
\end{sol}

\begin{exercise}[4a2]
    Prove that if $w, z \in \cbb$, then $||w| - |z|| \leq |w - z|$.
    
    \note{The inequality above is called the \emph{\textbf{reverse triangle inequality}}.}
\end{exercise}

\begin{sol}
    Let $w, z \in \cbb$.
	By the triangle inequality, we have
    \[|w| = |(w - z) + z| \leq |w - z| + |z|.\]
    Rearranging gives $|w| - |z| \leq |w - z|$.
	Similarly,
    \[|z| = |(z - w) + w| \leq |z - w| + |w| = |w - z| + |w|,\]
    which gives $|z| - |w| \leq |w - z|$, or equivalently $-|w - z| \leq |w| - |z|$.
	Combining these inequalities, we get
    \[||w| - |z|| \leq |w - z|. \qh\]
\end{sol}

\begin{exercise}[4a3]
    Suppose $V$ is a complex vector space and $\phi \in V'$.
	Define $\sigma: V \to \rbb$ by $\sigma(v) = \Re \phi(v)$ for each $v \in V$.
	Show that
    \[\phi(v) = \sigma(v) - i\sigma(iv)\]
    for all $v \in V$.
\end{exercise}

\begin{sol}
    Let $v \in V$.
	Then there exist $s, t \in \rbb$ such that
    \[\phi(v) = s + ti.\]
    Then
    \[\sigma(v) = \Re \phi(v) = s \quad \text{and} \quad \sigma(iv) = \Re \phi(iv) = \Re (i(s + ti)) = \Re (-t + si) = -t.\]
    Therefore,
    \[\sigma(v) - i\sigma(iv) = s - i(-t) = s + ti = \phi(v). \qh\]
\end{sol}

\begin{exercise}[4a4]
    Suppose $m$ is a positive integer.
	Is the set
    \[\set{0} \cup \set{p \in \pcal(\fbb) \mid \deg p = m}\]
    a subspace of $\pcal(\fbb)$?
\end{exercise}

\begin{sol}
    No, the set is not a subspace of $\pcal(\fbb)$.
	To see this, consider two polynomials $p, q \in \pcal(\fbb)$ with $p(x) = x^m$ and $q(x) = -x^m + x^{m - 1}$.
	Then $p + q = x^{m - 1}$, which has degree $m - 1$ and is not in the set.
	Therefore, the set is not closed under addition and hence not a subspace.
\end{sol}

\begin{exercise}[4a5]
    Is the set
    \[\set{0} \cup \set{p \in \pcal(\fbb) \mid \deg p \text{ is even}}\]
    a subspace of $\pcal(\fbb)$?
\end{exercise}

\begin{sol}
    No, the set is not a subspace of $\pcal(\fbb)$.
	To see this, consider two polynomials $p, q \in \pcal(\fbb)$ with $p(x) = x^2$ and $q(x) = x - x^2$.
	Then $p + q = x$, which has degree $1$ and is not even.
	Therefore, the set is not closed under addition and hence not a subspace.
\end{sol}

\begin{exercise}[4a6]
    Suppose that $m$ and $n$ are positive integers with $m \leq n$, and suppose $\lambda_1, \ldots, \lambda_m \in \fbb$.
	Prove that there exists a polynomial $p \in \pcal(\fbb)$ with $\deg p = n$ such that $0 = p(\lambda_1) = \cdots = p(\lambda_m)$ and such that $p$ has no other zeros.
\end{exercise}

\begin{sol}
    Consider the polynomial
    \[p(x) = (x - \lambda_1)^{n - m + 1}(x - \lambda_2) \cdots (x - \lambda_m).\]
    This polynomial has degree $n$, vanishes at $\lambda_1, \ldots, \lambda_m$, and has no other zeros.
	Therefore, it satisfies the required conditions.
\end{sol}

\begin{exercise}[4a7]
    Suppose that $m$ is a nonnegative integer, $z_1, \ldots, z_{m + 1}$ are distinct elements of $\fbb$, and $w_1, \ldots, w_{m + 1} \in \fbb$.
	Prove that there exists a unique polynomial $p \in \pcal_m(\fbb)$ such that
    \[p(z_k) = w_k\]
    for each $k = 1, \ldots, m + 1$.
    
    \note{This result can be proved without using linear algebra.
	However, try to find the clearer, shorter proof that uses some linear algebra.}
\end{exercise}

\begin{sol}
    Consider the linear map $T: \pcal_m(\fbb) \to \fbb^{m + 1}$ defined by
    \[T(p) = (p(z_1), \ldots, p(z_{m + 1})).\]
    Since $\dim \pcal_m(\fbb) = m + 1$ and $\dim \fbb^{m + 1} = m + 1$, it suffices to show that $T$ is injective to conclude that it is also surjective.
	Suppose $T(p) = 0$, i.e., $p(z_k) = 0$ for all $k = 1, \ldots, m + 1$.
	Since $p$ has degree at most $m$ and has $m + 1$ distinct zeros, it must be the zero polynomial.
	Therefore, $T$ is injective, and hence bijective.
	This shows that for any $(w_1, \ldots, w_{m + 1}) \in \fbb^{m + 1}$, there exists a unique $p \in \pcal_m(\fbb)$ such that $T(p) = (w_1, \ldots, w_{m + 1})$, which is exactly the desired polynomial.
\end{sol}

\begin{exercise}[4a8]
    Suppose $p \in \pcal(\cbb)$ has degree $m$.
	Prove that $p$ has $m$ distinct zeros if and only if $p$ and its derivative $p'$ have no zeros in common.
\end{exercise}

\begin{sol}
    Suppose $p$ has $m$ distinct zeros, say $z_1, \ldots, z_m$.
	If $p$ and $p'$ have a common zero $z_k$, then $z_k$ would be a multiple zero of $p$, contradicting the assumption that all zeros are distinct.
	Hence, $p$ and $p'$ have no zeros in common.

    Conversely, suppose $p$ and $p'$ have no zeros in common.
	If $p$ had a multiple zero $z$, then $p(z) = 0$ and $p'(z) = 0$, which would contradict the assumption.
	Therefore, all zeros of $p$ are distinct.
\end{sol}

\begin{exercise}[4a9]
    Prove that every polynomial of odd degree with real coefficients has a real zero.

    \apnote{4a9}{ provides an alternate solution using the intermediate value theorem.}
\end{exercise}

\begin{sol}
    Let $p(x)$ be a polynomial of odd degree with real coefficients.
	Since complex zeros come in pairs, there are an even number of complex zeros.
	Therefore, there must be at least one real zero.
\end{sol}

\begin{exercise}[4a10]
    For $p \in \pcal(\rbb)$, define $Tp: \rbb \to \rbb$ by
    \[(Tp)(x) = \begin{cases}
        \frac{p(x) - p(3)}{x - 3} & \text{if } x \neq 3, \\
        p'(3) & \text{if } x = 3
    \end{cases}\]
    for each $x \in \rbb$.
	Show that $Tp \in \pcal(\rbb)$ for every polynomial $p \in \pcal(\rbb)$ and also show that $T: \pcal(\rbb) \to \pcal(\rbb)$ is a linear map.
\end{exercise}

\begin{sol}
    Let $p \in \pcal(\rbb)$.
	Since $p(x) - p(3)$ has a zero at $x = 3$, we can factor it as $(x - 3)q(x)$ for some polynomial $q \in \pcal(\rbb)$.
	Then for $x \neq 3$,
    \[(Tp)(x) = \frac{p(x) - p(3)}{x - 3} = q(x),\]
    and differentiating $p$ gives $p'(3) = q(3)$.
	Since $Tp(3) = p'(3) = q(3)$, we have $Tp(x) = q(x)$ for all $x \in \rbb$.
	Therefore, $Tp \in \pcal(\rbb)$.

    Let $p_1, p_2 \in \pcal(\rbb)$ and $\alpha, \beta \in \rbb$.
	Let $p_1(x) - p_1(3) = (x - 3)q_1(x)$ and $p_2(x) - p_2(3) = (x - 3)q_2(x)$ for some $q_1, q_2 \in \pcal(\rbb)$.
	Then
    \[(\alpha p_1 + \beta p_2)(x) = (\alpha p_1(3) + \beta p_2(3)) + (x - 3)(\alpha q_1(x) + \beta q_2(x)).\]
    Hence, $T(\alpha p_1 + \beta p_2)(x) = \alpha q_1(x) + \beta q_2(x) = \alpha Tp_1(x) + \beta Tp_2(x)$ for all $x \in \rbb$.
	Therefore, $T$ is a linear map.
\end{sol}

\begin{exercise}[4a11]
    Suppose $p \in \pcal(\cbb)$.
	Define $q: \cbb \to \cbb$ by
    \[q(z) = p(z) \bar{p(\bar{z})}.\]
    Prove that $q$ is a polynomial with real coefficients.

    \apnote{4a11}{ provides an alternate solution by looking at the roots of $p$.}
\end{exercise}

\begin{sol}
    Let $p \in \pcal(\cbb)$ have the form $p(z) = a_0 + a_1 z + \cdots + a_n z^n$ with $a_k \in \cbb$.
	Then
    \[\bar{p(\bar z)} = \bar{\sum_{i = 0}^n a_i \bar z^i} = \sum_{i = 0}^n \bar a_i z^i.\]
    Therefore,
    \[q(z) = p(z) \bar{p(\bar z)} = \left(\sum_{i = 0}^n a_i z^i\right)\left(\sum_{j = 0}^n \bar a_j z^j\right) = \sum_{k = 0}^{2n} \left(\sum_{i + j = k} a_i \bar a_j\right) z^k.\]
    Let $b_m$ denote the coefficient of $z^m$ in $q(z)$.
	Then
    \[b_m = \sum_{i + j = m} a_i \bar a_j = \sum_{i + j = m} \bar a_j a_i = \bar{b_m}.\]
    Hence, $b_m \in \rbb$ for each $m = 0, 1, \ldots, 2n$.
	Therefore, $q$ is a polynomial with real coefficients.
\end{sol}

\begin{exercise}[4a12]
    Suppose $m$ is a nonnegative integer and $p \in \pcal_m(\cbb)$ is such that there are distinct real numbers $x_0, x_1, \ldots, x_m$ with $p(x_k) \in \rbb$ for each $k = 0, 1, \ldots, m$.
	Prove that all coefficients of $p$ are real.
\end{exercise}

\begin{sol}
    Let $p(x_k) = y_k \in \rbb$ for each $k = 0, 1, \ldots, m$.
	By \exref{4a7} with $\fbb = \rbb$, there exists a polynomial $q \in \pcal_m(\rbb)$ that satisfies $q(x_k) = y_k$ for each $k = 0, 1, \ldots, m$.
	Note that this polynomial must be the same as the original $p \in \pcal_m(\cbb)$ because a polynomial of degree at most $m$ is uniquely determined by its values at $m+1$ distinct points.
	Therefore, all coefficients of $p$ are real.
\end{sol}

\begin{exercise}[4a13]
    Suppose $p \in \pcal(\fbb)$ with $p \neq 0$.
	Let $U = \set{pq \mid q \in \pcal(\fbb)}$.
    \begin{subexercises}
        \item Show that $\dim \pcal(\fbb) / U = \deg p$.
        \item Find a basis of $\pcal(\fbb) / U$.
    \end{subexercises}
\end{exercise}

\begin{sole}
    \item We first show that $U$ is a subspace of $\pcal(\fbb)$.
	Clearly, $0 \in U$ since $0 = p \cdot 0$.
	If $u_1, u_2 \in U$, then $u_1 = pq_1$ and $u_2 = pq_2$ for some $q_1, q_2 \in \pcal(\fbb)$.
	For any scalars $\alpha, \beta \in \fbb$, we have $\alpha u_1 + \beta u_2 = \alpha pq_1 + \beta pq_2 = p(\alpha q_1 + \beta q_2) \in U$.
	Therefore, $U$ is closed under addition and scalar multiplication, and hence $U$ is a subspace of $\pcal(\fbb)$.
    
    Now let $r \in \pcal(\fbb)$.
	By the division algorithm for polynomials, there exist unique polynomials $q_0 \in \pcal(\fbb)$ and $s \in \pcal(\fbb)$ with $\deg s < \deg p$ such that $r = pq_0 + s$.
	This shows that every polynomial in $\pcal(\fbb)$ can be written as an element of $U$ plus a polynomial of degree less than $\deg p$.
	Consider the map $T : \pcal_{\deg p - 1}(\fbb) \to \pcal(\fbb) / U$ defined by $T(s) = s + U$.
	This map is linear and surjective, and $\nul T$ is trivial since if $s + U = U$, then $s \in U$, which implies $s = pq$ for some $q \in \pcal(\fbb)$.
	But $\deg s < \deg p$, so $s$ must be zero.
	Therefore, $T$ is an isomorphism, and we conclude that $\dim \pcal(\fbb) / U = \dim \pcal_{\deg p - 1}(\fbb) = \deg p$.
    \item A basis of $\pcal(\fbb) / U$ is given by the images of the standard basis of $\pcal_{\deg p - 1}(\fbb)$ under $T$, namely $\set{1 + U, x + U, \ldots, x^{\deg p - 1} + U}$.
\end{sole}

\begin{exercise}[4a14]
    Suppose $p, q \in \pcal(\cbb)$ are nonconstant polynomials with no zeros in common.
	Let $m = \deg p$ and $n = \deg q$.
	Use linear algebra as outlined below in (a)--(c) to prove that there exist $r \in \pcal_{n-1}(\cbb)$ and $s \in \pcal_{m-1}(\cbb)$ such that
    \[rp + sq = 1.\]
    \begin{subexercises}
        \item Define $T: \pcal_{n-1}(\cbb) \times \pcal_{m-1}(\cbb) \to \pcal_{m+n-1}(\cbb)$ by
        \[T(r, s) = rp + sq.\]
        Show that the linear map $T$ is injective.
        \item Show that the linear map $T$ in (a) is surjective.
        \item Use (b) to conclude that there exist $r \in \pcal_{n-1}(\cbb)$ and $s \in \pcal_{m-1}(\cbb)$ such that $rp + sq = 1$.
    \end{subexercises}
    \remark{In part (a), one may also observe that a multiple of $p$ and $q$ must have $\deg p + \deg q$, but $\deg rp = \deg r + \deg p \leq n - 1 + m$.}
\end{exercise}

\begin{sole}
    \item To show that $T$ is injective, suppose that $T(r, s) = rp + sq = 0$ for some $(r, s) \in \pcal_{n-1}(\cbb) \times \pcal_{m-1}(\cbb)$.
	Then $rp = -sq$ so that $q$ divides $rp$.
	Since $p$ and $q$ have no zeros in common, $q$ must divide $r$.
	But $\deg r < \deg q$, so the only possibility is that $r = 0$.
	Then $-sq = 0$.
	Since $q$ is nonconstant, this implies that $s = 0$.
	Therefore, $\nul T = \set{(0, 0)}$, and hence $T$ is injective.
    \item To show that $T$ is surjective, note that $\dim(\pcal_{n-1}(\cbb) \times \pcal_{m-1}(\cbb)) = n + m = \dim \pcal_{m+n-1}(\cbb)$.
	Since $T$ is a linear map between vector spaces of the same finite dimension, and we have already shown that $T$ is injective, it follows that $T$ is also surjective.
    \item Since $T$ is surjective, for the polynomial $1 \in \pcal_{m+n-1}(\cbb)$, there exist $r \in \pcal_{n-1}(\cbb)$ and $s \in \pcal_{m-1}(\cbb)$ such that $T(r, s) = rp + sq = 1$.
\end{sole}